\documentclass[11pt]{article}
\usepackage[a4paper,margin=23mm]{geometry}
\usepackage[T1]{fontenc}\usepackage{lmodern,microtype}
\usepackage{amsmath,amssymb,amsthm,booktabs,longtable,array,hyperref,fancyhdr}
\hypersetup{colorlinks=true,linkcolor=blue,urlcolor=blue}
\pagestyle{fancy}\fancyhf{}\fancyhead[L]{Rank 12, multiplicity one}\fancyhead[R]{Audit and continuation}\fancyfoot[C]{\thepage}
\setlength{\headheight}{14pt}\setlength{\parindent}{0pt}\setlength{\parskip}{5pt}
\widowpenalty=10000\clubpenalty=10000
\newtheorem{theorem}{Theorem}[section]\newtheorem{lemma}[theorem]{Lemma}\newtheorem{proposition}[theorem]{Proposition}
\newcommand{\FPdim}{\operatorname{FPdim}}\renewcommand{\Vec}{\operatorname{Vec}}\newcommand{\Rep}{\operatorname{Rep}}\newcommand{\Irr}{\operatorname{Irr}}
\newcommand{\ring}[1]{\texttt{R12\_1\_#1}}
\title{Rank-twelve multiplicity-one fusion rings\\\large Audit and theoretical continuation}
\author{Offline research record}\date{25 October 2026}
\begin{document}\maketitle
\begin{center}\small AI-generated research; independent human verification remains pending.\end{center}

\section{Results and audit conclusion}
For the same 632 supplied based rings, the updated decisions are
\[
 \boxed{632=@EXCLUDED@\ \text{excluded}+@MODELED@\ \text{unitarily modeled}+@UNRESOLVED@\ \text{unresolved}.}
\]
The first round recorded 393 exclusions, 127 unitary realizations and 112 unresolved
rings. The mathematical and implementation audit found no erroneous first-round
decision requiring retraction. The continuation resolves @NEWNEG@ further rings
by general complex obstructions and @NEWPOS@ by actual unitary constructions.
These counts concern based rings admitting categorifications, not tensor-equivalence
classes of categories. No repository or website has been revised.

The original input order and IDs are preserved. The decompressed input SHA-256 is
\begin{center}\scriptsize\texttt{6949e6fb8afd23f80b54e5a546a0dd34257430fd26c26b7acd8ff8444252f4d7}.\end{center}
The original verification of the fusion axioms and pairwise nonisomorphism is retained.
The present evidence includes every one of the 632 updated decisions, complete
unresolved tensors, construction data, exclusion certificates and independent audits.

\subsection*{What was checked in the first round}
An independent exclusion implementation, importing none of the first-round criterion
implementations, recovered exactly the 393 excluded IDs. It recomputed the arithmetic
and near-group premises, all 208 cyclic witnesses, the 12 field/order contradictions,
the six integral primary forms, all 4012 closed subring embeddings and the inherited
negative proofs. In particular, the Liptrap negatives cancel arbitrary nonzero complex
variables for every nondegenerate bicharacter assignment; failure of a finite phase
search was not used as a nonexistence argument. The rank-eight signed descent was
also checked with all 8192 unrestricted sign assignments, before indicator normalization.

The positive audit independently bound all 381 ordinary full-tensor identifications
and the 21 Liptrap reconstructions, accounting for all 402 construction records for
127 rings. It examined the categorical premises: actual tensor actions, central lifts,
zesting coherence, positive Fourier blocks and balanced bimodule products. This is
an audit of the recorded proofs and exact evidence, not a formalization of the cited
literature. Normal and assertion-disabled replays pass.

A reference-version issue was checked explicitly: the published two-dimension
reduction \cite[Lemma 3.1]{Sch} gives pseudo-unitarity \emph{after Galois conjugation}.
Literal pseudo-unitarity of every original category would be false (Yang--Lee is
a control). The first round already used the published, correct formulation, so no
classification change arose from this issue. Its restricted zest sampler was incomplete
as a discovery method, but its accepted models remain valid.

\section{A subgroup integrality criterion}
The following formulation covers noncyclic stabilizers without assuming that
independently chosen projective stabilizer cochains coincide.

\begin{theorem}\label{thm:stabilizer}
Let $\mathcal C$ be a complex fusion category, let $X,Y,Z$ be simple objects, and
let $H$ be a finite subgroup of its invertible objects. Suppose
\[
 hX\cong X,\qquad Yh\cong Y,\qquad hZ\cong Zh\cong Z\quad(h\in H),
 \qquad m=[Z:X\otimes Y]>0.
\]
Then $m\FPdim(Z)/|H|$ is an algebraic integer.
\end{theorem}
\begin{proof}
Choose isomorphisms $hX\to X$. Since $X$ is simple, their composition constants
define a connected twisted subgroup algebra $A$ with underlying object
$\bigoplus_{h\in H}h$, making $X$ a left $A$-module. The module pentagon supplies
the associativity equation, including the ambient pointed associator. Similarly,
$Y$ supplies a possibly different subgroup algebra $B$ and a right $B$-module.
These algebras are separable in characteristic zero and both have FP dimension $|H|$.

The commuting outer actions make $X\otimes Y$ an $A$--$B$ bimodule. Because $H$
fixes $Z$ on both sides, its full isotypic summand $M=mZ$ is stable under both actions.
Taking its Morita dual and the balanced product gives an object of
$\operatorname{Bimod}_{\mathcal C}(A)$ of FP dimension
\[
 \frac{\FPdim_{\mathcal C}(M\otimes_B M^\vee)}{\FPdim(A)}
 =\frac{m^2\FPdim(Z)^2}{\FPdim(A)\FPdim(B)}
 =\left(\frac{m\FPdim(Z)}{|H|}\right)^2.
\]
We use the standard relative dimension formula and bimodule normalization
\cite[Chapter 7]{EGNO}; see also \cite[Corollary 3.32 and Proposition 3.34]{Sanford}.
Thus this square is an algebraic integer. The number itself is integral over the
algebraic integers, so transitivity of integrality proves the claim.
\end{proof}

For the paired product $X\otimes X^*$ one can use the same algebra on both sides:
the transpose action makes $X^*$ a right $A$-module, and $mZ$ itself has normalized
FP dimension $m\FPdim(Z)/|H|$. Neither proof needs a pivotal, spherical or unitary
structure. The two-algebra argument explains why the general statement does not
require identifying independent noncyclic stabilizer cocycles.

The complete scan of all 632 supplied rings uses the intersection
$\operatorname{Stab}_L(X)\cap\operatorname{Stab}_R(Y)\cap
\operatorname{Stab}_L(Z)\cap\operatorname{Stab}_R(Z)$.
It conflicts with no accepted model and adds precisely the following sixteen exclusions.

\begin{proposition}
The sixteen rings with suffixes
\[
32,131,152,184,217,245,286,290,297,302,339,411,416,458,465,527
\]
have no complex fusion categorification.
\end{proposition}
\begin{proof}
Put $a=2\cos(\pi/7)$ and $b=a^2-1$. The exact dimension character is
$1^{\times8},(2a)^{\times2},(2b)^{\times2}$, with
$a^3-a^2-2a+1=0$. The certificate selects a simple $X$ of dimension $2b$ whose
order-four left stabilizer also fixes it on the right, and with $[X:X X^*]=1$.
Theorem~\ref{thm:stabilizer} would make $b/2$ integral. Its monic minimal polynomial is
\[
 t^3-t^2-\tfrac14t+\tfrac18,
 \qquad N_{\mathbb Q(a)/\mathbb Q}(b/2)=-\tfrac18.
\]
The nonintegral rational norm is impossible for an algebraic integer.
\end{proof}

\section{Morita enlargement excludes four more cubic rings}
We first prove a small lemma about every category with the indicated fusion rules,
rather than imposing equivalence to a particular known model.

\begin{lemma}\label{lem:s3}
Suppose $\mathcal B$ has simples $1,g,T$ with
$g^2=1$, $gT=Tg=T$ and $T^2=1+g+T$. For the subgroup algebra $A=1+g$
supplied by $T$, the dual category $\operatorname{Bimod}_{\mathcal B}(A)$ is pointed
of dimension six.
\end{lemma}
\begin{proof}
The left $A$-module category has exactly three simples $M_0=A,M_+,M_-$, all with
underlying dimension two. The last two are the two normalized $g$-actions on $T$.
The two invertible $A$-bimodules supported on $1+g$ are $A$ and $u=A_\sigma$, where
$\sigma$ changes the sign on $g$. The action of $u$ fixes $M_0$ and swaps $M_+,M_-$.

Every other simple bimodule is supported on $qT$ and has dimension $q$ in the dual.
Since the total dimension is six, the sole nonpointed alternative is a unique
dimension-two simple $Q$. It satisfies $Q^*=Q$ and $uQ=Q$. Therefore
$QM_0=M_++M_-$. Reciprocity makes the coefficient of $M_0$ in $QM_+$ exactly one.
All three module dimensions agree, so $QM_+$ has total multiplicity two. Its
$u$-invariance would require equal multiplicities of $M_+,M_-$. The remaining
odd total multiplicity cannot be distributed between that pair, a contradiction.
\end{proof}

\begin{proposition}
The four rings $\ring{38},\ring{39},\ring{218},\ring{219}$ have no complex
fusion categorification.
\end{proposition}
\begin{proof}
Put $b=2\cos(\pi/9)$, $x=1+2b$ and $y=2b^2-2$, so $b^3-3b-1=0$.
The source type is $1^{\times4},2^{\times2},x^{\times4},y^{\times2}$.
Its dimension-twelve integral subcategory $\mathcal B_0$ contains the subcategory
of Lemma~\ref{lem:s3}, with labels
\begin{center}\begin{tabular}{lrrr}\toprule Ring suffix&$g$&$T$&$h$ outside $\mathcal B$\\\midrule
38,39&11&10&1\\218&7&11&3\\219&6&11&3\\\bottomrule\end{tabular}\end{center}
Here $h$ is invertible and normalizes $\langle g\rangle$. Its conjugate subgroup
algebra is isomorphic to $A$: two multiplications on $1+g$ differ by rescaling $g$.
Thus $hA$ has an invertible bimodule structure outside the pointed dimension-six
dual. Consequently $\operatorname{Bimod}_{\mathcal B_0}(A)$ is pointed of order twelve.

The six nonintegral source simples are paired freely by $g$, with equal left/right
label actions. Each pair gives two simple bimodules, since its diagonal $C_2$
stabilizer has only ordinary one-dimensional projective representations over
$\mathbb C$. Thus the full dual has type $1^{\times12},x^{\times4},y^{\times2}$.
No multiplicity-one assertion about this dual is needed.

The numbers $1,x,y$ are rationally independent and $y^2=6+2x+2y$. For a
$y$-dimensional simple $Y$, its left and right stabilizers therefore both have
order six inside the twelve invertibles. Their intersection contains a subgroup
$H$ of order three. Each stabilizer fixes both $y$-dimensional labels, because
their orbits have at most two elements. Some $y$-dimensional $Z$ occurs in $YY^*$
with multiplicity $m=1$ or $2$. Theorem~\ref{thm:stabilizer} requires $my/3$ to be
integral, whereas
\[
 N_{\mathbb Q(b)/\mathbb Q}(y)=-24,\qquad
 N_{\mathbb Q(b)/\mathbb Q}(my/3)=-8m^3/9\notin\mathbb Z.
\]
This contradicts integrality in both cases.
\end{proof}

\section{Four quadratic exclusions via a linear sign reduction}
The rings $\ring{63},\ring{64},\ring{215},\ring{365}$ have exact dimensions
\[
 1^{\times2},\quad d^{\times8},\quad(d+1)^{\times2},\qquad
 d=4+3\sqrt2,\quad d^2=8d+2.
\]
Their common Casimir polynomial is
\[
 (t-36)(t-18)^8(t-4)(t^2-720t+2592).
\]
We justify sphericality before using any spherical character obstruction.

\begin{lemma}\label{lem:parity}
In a pivotal category, a self-dual simple with second indicator $-1$ has even
multiplicity in every $Y\otimes Y^*$.
\end{lemma}
\begin{proof}
The object $M=Y\otimes Y^*$ has the symmetric self-duality
$J=j_Y\otimes\mathrm{id}_{Y^*}:M\to M^*$, where $j$ is the pivotal structure.
The identity $J^*j_M=J$ follows from $j_{Y^*}=(j_Y^*)^{-1}$.
On a self-dual simple's isotypic part, write $J=\phi\otimes B$. Transposition gives
$B^{\mathsf T}=\nu_2(\phi)B$. For indicator $-1$, the nondegenerate multiplicity
pairing is alternating, so its dimension is even.
\end{proof}

Choose dual-compatible lifts in the canonical sphericalization and the
positive-second-indicator lift for each self-dual simple. Such a simple has
self-dual lifts: exchanging them by duality would equate opposite nonzero spherical
dimensions. Tensoring with the central sign object changes the second indicator's sign.
Writing the signed quotient coefficient as $T=a_+-a_-$, Lemma~\ref{lem:parity}
gives $N_{Y,Y^*}^S-T_{Y,Y^*}^S\equiv0\pmod4$ for the selected positive lift $S$.

Every nonzero coefficient is $(-1)^v$ for a bit $v$. Cyclic trace and dual reversal
identify its variable orbit. Associativity modulo four is linear over $\mathbb F_2$:
in each nonzero associativity term, add the two coefficient bits, and equate the
sum on both sides. Ordinary associativity already equates the term counts.
Together with the preceding parity conditions this gives the following short,
complete reduction:
\begin{center}\begin{tabular}{rrrrr}\toprule Suffix&Variables&Kernel dimension&Linear choices&Associative choices\\\midrule
63&189&1&2&2\\64&190&1&2&2\\215&179&2&4&4\\365&169&4&16&8\\\bottomrule\end{tabular}\end{center}
Every surviving signed tensor is a change of simple basis signs of $N$, with
unit sign one and equal signs on dual pairs. The exact gauges are recorded.
Changing the selected lifts accordingly makes all their products positive.
These lifts form a full spherical tensor subcategory equivalent to the original
category. Thus any hypothetical categorification is spherical. This uses a
rank-independent categorical lemma and new small linear systems; the frozen
rank-eight engine is not silently enlarged.

We compute all one-dimensional characters, including in the noncommutative cases.
The common annihilator of the regular commutators is exactly the sum of the
one-dimensional simple blocks. Restricting the fusion matrices there and choosing
a square-free separating element gives every character orbit by exact polynomial
coordinates. An independent audit instead uses the commutator-ideal quotient.

A spherical character has no zero simple dimensions and takes equal values on
dual pairs. Apart from its FP quadratic orbit, every surviving character field has
degree one or three. But all formal codegrees must lie in that field: the
distinguished central simple has dimension $\dim(\mathcal C)/f_E$ and forgetting it
is an integer sum of source simple dimensions \cite{Ostrik}. The Casimir quadratic
$t^2-720t+2592$ cannot split in a field of odd degree. These branches are impossible.

In the remaining quadratic orbit the dimension order is $\mathbb Z+\mathbb Z d$
and $D=24+84d$. The Casimir root 36 has multiplicity one, so its irreducible
representation degree is one and it is the actual formal codegree 36. Its central
simple would have dimension
\[
 D/36=\tfrac23+\tfrac73d=10+7\sqrt2\notin\mathbb Z+\mathbb Z d.
\]
The same order coordinates are nonintegral in the other embedding. This excludes
every spherical character, and therefore every complex categorification of the four rings.

@D40SECTION@

\section{Additional unitary models}
The following table lists only accepted categorical constructions, each matched
against all 1728 input coefficients. A proposed action or algebra whose existence
has not been established is not counted.
@MODELTABLE@

\subsection*{The graph-subgroup model at dimension 120}
Let $G=A_5\times C_2$, let $D_{10}\leq A_5$ be generated by a five-cycle and a
reflection, and let $\chi:D_{10}\to C_2$ be the reflection-sign character. The
actual subgroup $H=\{(h,\chi(h)):h\in D_{10}\}$ has order ten. The connected
separable subgroup algebra in $\Vec_G$ gives the unitary category
$\mathcal C(G,1,H,1)$, realizing $\ring{25}$.
Its double cosets and stabilizer characters give four invertibles, four
dimension-two simples and four dimension-five simples. Direct convolution and
an independent construction of all induced graded bimodules and balanced products
give the same tensor. Unlike $D_{10}\subset S_5$, this graph subgroup surjects
onto the extra $C_2$ factor; the unwanted quotient grading is absent.

\subsection*{A concrete twisted binary-tetrahedral model}
The category realizing $\ring{597}$ is $\mathcal C(SL(2,3),\omega,C_4,\psi)$.
Here is a source-independent specification of the cocycle. Use
$SL(2,3)/\{\pm1\}\cong V_4\rtimes C_3$. For $v,w,u\in\mathbb F_2^2$ put
\[
 f(v,w,u)=v_0w_0u_0+v_1w_1u_1+v_0w_1u_1,
 \qquad \bar f=\sum_{j=0}^2 f(A^jv,A^jw,A^ju),
\]
where $A$ is the order-three action. This is an invariant trilinear cocycle.
On quotient coordinates $(v,a),(w,b),(u,c)$ define
\[
 \omega=(-1)^{\bar f(v,A^a w,A^{a+b}u)}
\]
and inflate it to $SL(2,3)$. The stored unit-modulus cochain $\psi$ on the actual
cyclic subgroup of order four satisfies its algebra equation exactly. Thus the
algebra exists without an asserted abstract classification of $H^3(SL(2,3))$.
Every group cocycle equation, subgroup equation, bimodule action and balanced
tensor coefficient is checked. This is a positive construction; it makes no
completeness claim about all group cocycles.

\subsection*{Triality, equivariantization and Morita constructions}
The full $S_3$ gaugings of the three-fermion category $SO(8)_1$ give
$\ring{79}$ and $\ring{80}$ \cite[Section 6.2]{CGPW}\cite[Section VI.D]{THF}.
These are actual unitary gauging constructions; the modular matrices are used
to identify their resulting fusion rules, not to infer existence from modular data.
Their sign-representation boson has an actual central lift and lies in the neutral
component of the integral/nonintegral $C_2$ grading. Its associative zest gives
$\ring{322}$ and $\ring{323}$.

For $\ring{449}$, inversion of $C_5$ preserves the bicharacter
$\chi(a,b)=\zeta_5^{ab}$ of a unitary TY category. It defines a weak $C_4$ action
with carry coherence $\kappa^{\lfloor(a+b)/4\rfloor}$, where $\kappa$ is one on
pointed objects and $-1$ on the TY-odd simple. The strict inversion functor has
order two; it is not claimed to have order four on labels. The invariant carry
coherence supplies the genuine unitary $C_4$ action. Its equivariantization has
type $1^{\times4},2^{\times4},(\sqrt5)^{\times4}$ and the full target tensor.

The additional dimension-40 Morita models use a unitary Liptrap reconstruction
over $C_5$ \cite[Theorem 4 and Lemma 5.4]{Liptrap}, optionally twisted by a stored
unit-modulus grading cocycle. The required order-two subgroup algebra and all
associators entering its bimodule products are supplied explicitly. Independent
checks evaluate the scalar pentagons, Fourier unitarity and balanced Hom spaces.

Finally, $\ring{541}$ is classical: $\Rep(C_3\rtimes C_8)$, with inversion action.
There are eight linear characters $r_k$ and four induced degree-two characters $X_j$,
with
\[
 r_kr_l=r_{k+l},\qquad r_kX_j=X_{j+k},\qquad
 X_iX_j=r_{i+j}+r_{i+j+4}+X_{i+j}.
\]
Subscripts on $r$ and $X$ are taken modulo eight and four, respectively.
The source-to-model permutation is $(0,2,6,1,7,9,11,5,3,8,10,4)$.

\section{Remaining cases and precise limits}
@REMAININGTABLE@
All @UNRESOLVED@ remaining IDs and tensors are supplied separately. Every case
received structural analysis and relevant model/criterion attempts. This remains
a partial categorification result, not a completed rank-twelve classification.

The dimension-48 central-quotient calculations identify compatible smaller rings
among the supplied data. They do not prove uniqueness of their categorifications
or tensor actions, and were not used as negative completeness claims. The
four dimension-144 rings with suffixes 180,181,320,321 admit exact graded twists
at the based-ring level, but the necessary coherent tensor actions have not been
established. They remain unresolved.

For the five two-dimension cases with suffixes 332,350,561,623,628, the short primary
probes gave no negative certificate. In case 332, the trivial outer automorphism
group of a familiar unitary Haagerup model does not classify arbitrary original
categorifications; twisted pointed associators must also be covered. In case 350,
a conditional subgroup-dual reduction cannot be applied until its algebra exists
in every relevant pointed associator class. Similar central-lift and algebra
existence gaps are stated in the research notes. No such gap has been hidden by
a successful scalar calculation or a failed construction search.

Known and proposed constructions are distinguished throughout. ``New'' here concerns
the recorded argument or explicit construction, when it has not to our knowledge
been explicitly recorded in the literature; no novelty up to tensor equivalence
or exoticness is asserted. Independent human verification remains pending.

\section{Reproducibility}
The complete offline package preserves the first-round evidence and adds the
independent audit, all new accepted proofs and constructions, and the updated
JSON/CSV ledger. The top-level replay checks source bindings and every new
decision; the independent mathematical audits and constructor-specific replays
are included. Normal and optimized Python verification pass. Short exhaustive
linear/cohomological systems are explicitly distinguished from unsuccessful or
incomplete discovery attempts. No long frame search or numerical pentagon solving
was used. Verification of identities is not presented as a substitute for the
categorical existence and scope arguments above.

\begin{thebibliography}{99}\footnotesize\setlength{\itemsep}{2pt}\setlength{\parskip}{0pt}
\bibitem{EGNO} P.~Etingof, S.~Gelaki, D.~Nikshych and V.~Ostrik,
\emph{Tensor Categories}, AMS, 2015, Chapter 7.
\bibitem{Ostrik} V.~Ostrik, \emph{On formal codegrees of fusion categories},
Math. Res. Lett. 16 (2009), 895--901; \href{https://arxiv.org/abs/0810.3242}{arXiv:0810.3242}.
\bibitem{Sanford} S.~Sanford, \emph{Fusion Categories over Non-Algebraically Closed Fields},
\href{https://arxiv.org/abs/2401.02354}{arXiv:2401.02354v3}, Corollary 3.32 and Proposition 3.34.
\bibitem{Sch} A.~Schopieray, \emph{Categorification of integral group rings extended by one dimension},
J. London Math. Soc. 108 (2023), 1617--1641; \href{https://arxiv.org/abs/2208.07319}{arXiv:2208.07319}. Published Lemma 3.1.
\bibitem{ENO} P.~Etingof, D.~Nikshych and V.~Ostrik,
\emph{Fusion categories and homotopy theory}, Quantum Topol. 1 (2010), 209--273;
\href{https://arxiv.org/abs/0909.3140}{arXiv:0909.3140}, extension classification and $\operatorname{BrPic}(\Vec_A)$.
\bibitem{EGO} P.~Etingof, S.~Gelaki and V.~Ostrik,
\emph{Classification of fusion categories of dimension $pq$},
\href{https://arxiv.org/abs/math/0304194}{arXiv:math/0304194v2}, Lemmas 6.1--6.2.
\bibitem{GP} C.~Galindo and J.~Y.~Plavnik,
\emph{Tensor functors between Morita duals of fusion categories},
\href{https://arxiv.org/abs/1407.2783}{arXiv:1407.2783v4}, Theorem 7.9 and Corollary 7.11.
\bibitem{Liptrap} J.~Liptrap, \emph{Generalized Tambara--Yamagami categories},
\href{https://arxiv.org/abs/1002.3166v2}{arXiv:1002.3166v2}, Definition 4.3, Theorem 4 and Lemma 5.4.
\bibitem{CGPW} S.~X.~Cui, C.~Galindo, J.~Y.~Plavnik and Z.~Wang,
\emph{On Gauging Symmetry of Modular Categories},
\href{https://arxiv.org/abs/1510.03475}{arXiv:1510.03475v3}, Section 6.2.
\bibitem{THF} J.~C.~Y.~Teo, T.~L.~Hughes and E.~Fradkin,
\emph{Theory of Twist Liquids: Gauging an Anyonic Symmetry},
Ann. Phys. 360 (2015), 349--445; \href{https://arxiv.org/abs/1503.06812}{arXiv:1503.06812v2}, Section VI.D, Table IX and Eq. (6.53).
\bibitem{Zest} C.~Delaney, C.~Galindo, J.~Y.~Plavnik, E.~C.~Rowell and Q.~Zhang,
\emph{Braided zesting and its applications},
\href{https://arxiv.org/abs/2005.05544}{arXiv:2005.05544}, Section 3.
\end{thebibliography}
\end{document}
